\documentclass[11pt,a4paper]{article}

%% ── Packages ─────────────────────────────────────────────────────────────────
\usepackage[T1]{fontenc}
\usepackage[utf8]{inputenc}
\usepackage[margin=2.5cm]{geometry}
\usepackage{amsmath,amssymb,amsthm}
\usepackage{booktabs}
\usepackage{array}
\usepackage{hyperref}
\usepackage{parskip}

\hypersetup{
  colorlinks=true,
  linkcolor=blue,
  urlcolor=blue,
  citecolor=blue
}

%% ── Theorem environments ─────────────────────────────────────────────────────
\newtheorem{theorem}{Theorem}[section]

%% ── Title ────────────────────────────────────────────────────────────────────
\title{\textbf{Stage 2 Report}\\[0.4em]
  \large Feedback Linearization for Operational Space Trajectory Tracking\\
  of a Serial Manipulator (KUKA IIWA 14 R820)}
\author{Name: Sashish Jha\\Student ID: 27509}
\date{March 31, 2026}

\begin{document}

\maketitle
\thispagestyle{empty}

%% ─────────────────────────────────────────────────────────────────────────────
\section{Review of Existing Methodology}

\subsection{System Dynamics}

The KUKA IIWA 14 R820 is a 7-DOF serial manipulator whose joint-space dynamics
are governed by the standard rigid-body equation of motion:
\begin{equation}
  M(q)\ddot{q} + C(q,\dot{q})\dot{q} + G(q) = \tau,
  \label{eq:joint_dynamics}
\end{equation}
where $q \in \mathbb{R}^{7}$ is the joint position vector, $M(q)$ is the
symmetric positive-definite inertia matrix, $C(q,\dot{q})$ collects Coriolis
and centrifugal effects, $G(q)$ is the gravity vector, and $\tau$ is the
applied joint torque.

\subsection{Operational Space Feedback Linearization}

The operational-space (Cartesian-space) position of the end-effector is related
to joint coordinates through the forward kinematics map $X = f(q) \in
\mathbb{R}^{3}$. Differentiating twice yields:
\begin{equation}
  \ddot{X} = J(q)\ddot{q} + \dot{J}(q,\dot{q})\dot{q},
  \label{eq:xddot}
\end{equation}
where $J(q) = \partial f/\partial q \in \mathbb{R}^{3\times 7}$ is the
analytical Jacobian. Substituting the joint-space dynamics \eqref{eq:joint_dynamics}
into \eqref{eq:xddot} and choosing the torque input as:
\begin{equation}
  \tau = M(q)J^{-1}(q)\!\left[a_x - \dot{J}(q,\dot{q})\dot{q}\right]
       + C(q,\dot{q})\dot{q} + G(q),
  \label{eq:fl_law}
\end{equation}
the closed-loop Cartesian acceleration simplifies to:
\begin{equation}
  \ddot{X} = a_x,
  \label{eq:double_integrator}
\end{equation}
i.e.\ a pure double-integrator in Cartesian space. The auxiliary input $a_x$
is then chosen as a PD outer-loop controller:
\begin{equation}
  a_x = \ddot{X}_d + K_d\dot{e} + K_p e, \qquad e = X_d - X,
  \label{eq:pd_outer}
\end{equation}
where $K_p, K_d \in \mathbb{R}^{3\times 3}$ are positive-definite gain matrices.

\subsection{Limitations of the Baseline Method}

The control law \eqref{eq:fl_law} relies on exact cancellation of the nonlinear
terms $M(q)$, $C(q,\dot{q})$, and $G(q)$. In practice, three principal failure
modes arise.

\paragraph{Model uncertainty.}
Robot parameters (link masses, inertias, friction) are never perfectly known.
Any mismatch $\Delta M$, $\Delta C$, $\Delta G$ propagates directly as a
disturbance to the linearised system, degrading tracking accuracy.

\paragraph{Near-singular configurations.}
The term $J^{-1}(q)$ is ill-conditioned near kinematic singularities, causing
large torque spikes and potential instability.

\paragraph{No active disturbance rejection.}
External forces or payload variations are not accounted for; the baseline PD
law has no integral action and no explicit disturbance observer, so steady-state
errors can persist under constant perturbations.

%% ─────────────────────────────────────────────────────────────────────────────
\section{Proposed New Methodology}

\subsection{Motivation}

To address the limitations above, the proposed methodology augments the baseline
feedback-linearization controller with a \textbf{Sliding Mode Control (SMC)}
outer loop. SMC is chosen because: (i) it provides robust, finite-time
convergence to the desired trajectory even under bounded model uncertainty;
(ii) it requires only upper bounds on the uncertainty, which can be estimated
experimentally; and (iii) the design integrates naturally with the
double-integrator structure obtained after feedback linearization.

\subsection{Sliding Surface Design}

Define the tracking error in Cartesian space as $e = X_d - X \in \mathbb{R}^3$.
A first-order sliding surface is chosen as:
\begin{equation}
  s = \dot{e} + \Lambda e, \qquad \Lambda = \mathrm{diag}(\lambda_1, \lambda_2, \lambda_3),
  \quad \lambda_i > 0.
  \label{eq:sliding_surface}
\end{equation}
When $s = 0$ the error dynamics reduce to $\dot{e} + \Lambda e = 0$, which is
globally exponentially stable for any $\lambda_i > 0$.

\subsection{Control Law}

Recalling that the feedback-linearized system satisfies $\ddot{X} = a_x$, the
proposed auxiliary input is:
\begin{equation}
  a_x = \ddot{X}_d + \Lambda\dot{e} + K\,\mathrm{sgn}(s) + \eta s,
  \label{eq:smc_ax}
\end{equation}
where $K = \mathrm{diag}(k_1, k_2, k_3)$ with $k_i > \|\delta_i\|_\infty$ is
the switching gain ($\delta_i$ being the $i$-th component of the lumped
disturbance), and $\eta > 0$ provides additional exponential decay on the
sliding surface. Substituting \eqref{eq:smc_ax} into the linearised dynamics
gives:
\begin{equation}
  \dot{s} = -K\,\mathrm{sgn}(s) - \eta s + d(t),
  \label{eq:sdot}
\end{equation}
where $d(t)$ represents the lumped disturbance (model uncertainty plus external
forces). A standard Lyapunov argument with $V = \tfrac{1}{2}s^\top s$ yields:
\begin{equation}
  \dot{V} = s^\top\dot{s}
          \leq -\sum_{i=1}^{3}(k_i - \|d_i\|_\infty)|s_i| - \eta\|s\|^2 < 0,
  \label{eq:lyapunov}
\end{equation}
provided $k_i > \|d_i\|_\infty$. This guarantees finite-time reachability of
the sliding surface and subsequent exponential convergence to $e = 0$.

\subsection{Chattering Reduction}

The discontinuous $\mathrm{sgn}(\cdot)$ function in \eqref{eq:smc_ax} produces
high-frequency chattering that excites unmodelled joint flexibility. To mitigate
this, the standard saturation approximation is used:
\begin{equation}
  \mathrm{sgn}(s_i) \approx \mathrm{sat}\!\left(\frac{s_i}{\phi}\right)
  =
  \begin{cases}
    s_i/\phi, & |s_i| \leq \phi,\\
    \mathrm{sgn}(s_i), & |s_i| > \phi,
  \end{cases}
  \label{eq:sat}
\end{equation}
where $\phi > 0$ is a thin boundary-layer width. This replaces the ideal
sliding mode with a quasi-sliding mode, confining the error to an $O(\phi)$
neighbourhood of the surface.

\subsection{Complete Controller Summary}

The full proposed torque command is:
\begin{equation}
  \tau = M(q)J^{-1}(q)\!\left[a_x^{\mathrm{SMC}} - \dot{J}(q,\dot{q})\dot{q}\right]
       + C(q,\dot{q})\dot{q} + G(q),
  \label{eq:full_law}
\end{equation}
with $a_x^{\mathrm{SMC}}$ given by \eqref{eq:smc_ax}--\eqref{eq:sat}. The
inner feedback-linearization loop cancels the nominal nonlinear dynamics; the
outer SMC loop drives the Cartesian error to zero robustly.

%% ─────────────────────────────────────────────────────────────────────────────
\section{Analysis}

\subsection{Stability}

\begin{theorem}[Reachability]
Under the control law \eqref{eq:full_law}, if $k_i > \|d_i\|_\infty$ for each
axis $i$, the sliding surface $s = 0$ is reached in finite time
$t_r \leq \|s(0)\| / (k_{\min} - \|d\|_\infty)$.
\end{theorem}

\begin{theorem}[Tracking convergence]
Once on the surface, the error evolves as $\dot{e} = -\Lambda e$, yielding
$\|e(t)\| \leq \|e(t_r)\|e^{-\lambda_{\min}(t - t_r)}$.
\end{theorem}

These results hold even when the nominal plant parameters used to compute
$M(q)$, $C(q,\dot{q})$, $G(q)$ differ from the true values, as long as the
disturbance bound is satisfied.

\subsection{Comparison with Baseline}

\begin{center}
\begin{tabular}{@{}lcc@{}}
\toprule
\textbf{Property} & \textbf{Baseline FL + PD} & \textbf{Proposed FL + SMC} \\
\midrule
Nominal tracking               & Exponential          & Finite-time + exponential \\
Robustness to $\Delta M, \Delta G$ & None             & Guaranteed (bounded $d$)  \\
Steady-state error (const.\ dist.) & Non-zero         & Zero (with ideal sgn)     \\
Chattering                     & Absent               & Present (mitigated by $\phi$) \\
Gain tuning parameters         & $K_p, K_d$           & $\Lambda, K, \eta, \phi$  \\
\bottomrule
\end{tabular}
\end{center}

\subsection{Parameter Selection Guidelines}

The sliding surface bandwidth $\Lambda$ is chosen to match the desired
closed-loop bandwidth $\omega_c$ of the Cartesian axes: $\lambda_i \approx
\omega_{c,i}$. The switching gain $k_i$ must exceed the estimated worst-case
disturbance magnitude on each axis, which can be identified from open-loop
torque residuals. The boundary-layer width $\phi$ trades steady-state precision
against chattering amplitude and is typically set to 1--5\% of the nominal error
magnitude.

%% ─────────────────────────────────────────────────────────────────────────────
\section*{References}

\begin{enumerate}
  \item O.~Khatib, ``A unified approach for motion and force control of robot
    manipulators: The operational space formulation,'' \textit{IEEE Journal on
    Robotics and Automation}, vol.~3, no.~1, pp.~43--53, 1987.\\
    \url{https://ieeexplore.ieee.org/document/1087068}\\
    \textit{Relevance:} Establishes the operational space control framework that
    underpins the Cartesian-space trajectory tracking approach used in this
    project.

  \item M.~W.~Spong, ``Modeling and control of elastic joint robots,''
    \textit{IEEE Transactions on Robotics and Automation}, vol.~3, no.~6,
    pp.~291--300, 1987.\\
    \url{https://ieeexplore.ieee.org/document/1087035}\\
    \textit{Relevance:} Provides the theoretical basis for feedback linearization
    applied to robotic manipulators, which is the core control technique employed
    in this project.

  \item V.~I.~Utkin, \textit{Sliding Modes in Control and Optimization}.
    Springer, 1992.\\
    \textit{Relevance:} Foundational text on sliding mode control theory,
    providing the stability analysis framework adopted in Section~2.

  \item J.-J.~E.~Slotine and W.~Li, \textit{Applied Nonlinear Control}.
    Prentice Hall, 1991.\\
    \textit{Relevance:} Covers sliding mode design for robotic manipulators
    (Chapter~7), including the boundary-layer chattering reduction technique
    used in Section~2.4.
\end{enumerate}

\end{document}
